% ---- Outlook for habitability (joint synthesis) ----
\section{Outlook for habitability}

Taken together, the five studies make a layered rather than a singular case for Titan's habitability,
and point to one explicit conclusion: the most promising sites for prebiotic chemistry are the
liquid-water environments that impacts create. Impact melt pools -- and the subsurface ocean, if it
persists -- are the principal settings that bring liquid water into contact with Titan's abundant
surface organics \citep{Nicolaides2026, Crick2026}, and reactive simulations show that such a brine
spontaneously forms carbon--nitrogen bonds, mobilises amino acids, and incorporates phosphorus and
sulphur: the first steps of prebiotic chemistry \citep{Castlevetro2026}.

This potential is contingent, not guaranteed. Whether a given impact yields a viable reactor depends on
the ambient conditions when it occurs: the impactor's size, speed and angle set how much organic
feedstock survives atmospheric entry and shock, while the impact energy, target composition and epoch
set the melt volume and the pocket's freezing lifetime \citep{Nicolaides2026}. Habitability on Titan is
therefore episodic and site-specific rather than global.

Most speculatively, wherever liquid water endures -- the ocean, or long-lived melt lenses such as that
beneath Selk crater -- Titan could host anaerobic, methanogen-like life, whose biogenic methane is even
one candidate explanation for the moon's otherwise puzzling atmospheric supply \citep{IslamRahim2026}.
NASA's Dragonfly mission (2034) will sample impact-altered organics directly at Selk crater -- the
equatorial site our map singles out for its organic--ice contact terrain \citep{Lorenz2021,
Meadows2026} -- offering a first test; JWST can extend the
atmospheric record, while resolving the ocean question ultimately requires a dedicated gravity orbiter.
Should a red-giant Sun ultimately warm the surface, Titan's most habitable chapter may still lie ahead.
